\section*{Physique (13 points)}
\section*{Exercice 1 (3 points) : Datation par la méthode Uranium-Thorium}
Les sédiments marins contiennent du thorium ${ }_{90}^{230}$ Th et de l'uranium ${ }_{92}^{234}$ U avec des pourcentages différents selon leurs âges. Le thorium ${ }_{90}^{230}$ Th présent dans ces sédiments provient de la désintégration spontanée de l'uranium ${ }_{92}^{234} \mathrm{U}$ au cours du temps.\\
Le but de l'exercice est l'étude de la désintégration de l'uranium ${ }_{92}^{234}$ U .

\section*{Données :}
Energies de masse des nucléons et du noyau de l'uranium 234:

\begin{center}
\begin{tabular}{|c|c|c|c|}
\cline { 2 - 4 }
\multicolumn{1}{c|}{} & 92 protons & 142 neutrons & Noyau ${ }_{92}^{234} U$ \\
\hline
Énergie de masse en $(M e V)$ & 86321,9 & 133418,5 & 218009,1 \\
\hline
\end{tabular}
\end{center}

\begin{enumerate}
  \item Donner la composition du noyau de thorium ${ }_{90}^{230} T h$.
  \item Écrire l'équation de désintégration du noyau d'uranium ${ }_{92}^{234} U$. Identifier le type de cette désintégration.
  \item Recopier sur votre copie le numéro de la question et écrire la lettre correspondante à la proposition vraie.\\
L'énergie de liaison du noyau ${ }_{92}^{234} U$ vaut :
\end{enumerate}

\begin{center}
\begin{tabular}{|c|c|c|c|c|c|c|c|}
\hline
$\mathbf{A}$ & $1,65.10^{3} \mathrm{MeV}$ & $\mathbf{B}$ & $1,73.10^{3} \mathrm{MeV}$ & $\mathbf{C}$ & $1,85.10^{3} \mathrm{MeV}$ & $\mathbf{D}$ & $1,98.10^{3} \mathrm{MeV}$ \\
\hline
\end{tabular}
\end{center}

\begin{enumerate}
  \setcounter{enumi}{3}
  \item On considère un échantillon de sédiment marin qui s'est formé à l'instant $t_{0}=0$. Cet échantillon contient $N_{0}$ noyaux d'uranium et pas de noyaux de thorium.\\
On désigne par $a_{0}$ l'activité radioactive de l'échantillon à l'instant $t_{0}=0$ et par $a$ l'activité radioactive de l'échantillon à l'instant $t$.\\
La courbe ci-contre représente les variations de $\ln \left(\frac{a_{0}}{a}\right)$ en fonction du temps.\\
4.1. Déterminer graphiquement en unité ( $a n^{-1}$ ) la valeur de la constante radioactive $\lambda$ de l'uranium 234.\\
4.2. L'étude de l'échantillon à l'instant $t_{1}$ (âge de l'échantillon) a montré que $\frac{a_{0}}{a}=\sqrt{2}$.\\
	\includegraphics[width=.7\textwidth]{22adf374-91fd-4329-b77c-a241db44a97c-3_638_659_1626_1308}
	

\end{enumerate}

Déterminer en unité (an) la valeur de $t_{1}$ âge de l'échantillon.

\section*{Exercice 2 ( $\mathbf{5}$ points) : Étude de la réponse d'un dipôle}
Les circuits électriques ou électroniques comportent des condensateurs et des bobines dont les comportements diffèrent selon leurs usages.\\
Cet exercice vise :

\begin{itemize}
  \item l'étude de la réponse d'un dipôle RC à un échelon de tension ascendant;
  \item l'étude des oscillations électriques libres et l'échange énergétique dans un circuit RLC série.
\end{itemize}

On réalise le montage électrique représenté dans la figure (1) constitué des éléments suivants :

\begin{itemize}
  \item un générateur idéal de tension de force électromotrice $E$;
  \item un condensateur de capacité $C$ initialement non chargé;
  \item une bobine ( $L, r=0$ ) ;
  \item deux conducteurs ohmiques de résistances respectives $R_{1}=6 k \Omega$ et $R_{2}$;
  \item un interrupteur $K$.
\end{itemize}

\section*{1. Réponse d'un dipôle RC à un échelon de tension ascendant}
À l'instant $t_{0}=0$, on place l'interrupteur en position (1). La figure (2) représente la variation de la tension $u_{C}(t)$ aux bornes du condensateur.

\begin{figure}[H]
	\centering
	\includegraphics[width=.7\textwidth]{22adf374-91fd-4329-b77c-a241db44a97c-4_558_663_790_356}
	\caption{Figure (1)}
\end{figure}


\begin{figure}[H]
	\centering
	\includegraphics[width=.7\textwidth]{22adf374-91fd-4329-b77c-a241db44a97c-4_617_799_788_1043}
	\caption{Figure (2)}
\end{figure}


1.1. Montrer que l'équation différentielle vérifiée par $u_{C}$ s'écrit : $\frac{d u_{C}}{d t}+\frac{1}{\tau} \cdot u_{C}=\frac{E}{\tau}$ avec $\tau$ une constante positive. Donner l'expression de $\tau$.\\
1.2. Déterminer graphiquement les valeurs de $E$ et $\tau$.\\
1.3. Vérifier que $C \simeq 6,3 \mu F$.

\section*{2. Étude des oscillations électriques libres et échange énergétique}
Lorsque le régime permanent est atteint, on bascule l'interrupteur $K$ en position (2) à l'instant $t_{0}=0$.\\
La courbe de la figure (3) représente la variation de la tension $u_{C}(t)$ aux bornes du condensateur.\\
2.1. Justifier la nature des oscillations électriques dans le circuit.\\
2.2. Déterminer la valeur de la charge $Q_{0}$ du condensateur à l'instant $t_{0}=0$.

\begin{figure}[H]
	\centering
	\includegraphics[width=.7\textwidth]{22adf374-91fd-4329-b77c-a241db44a97c-4_508_1061_1827_904}
	\caption{Figure (3)}
\end{figure}


2.3. Déterminer graphiquement la valeur de la pseudo-période $T$ des oscillations.\\
2.4. En considérant que la pseudo-période $T$ est égale à la période propre de l'oscillateur ( $L C$ ), déterminer la valeur de l'inductance $L$ de la bobine (On prend $\pi^{2}=10$ ).\\
2.5. Les courbes de la figure (4) représentent les variations en fonction du temps de l'énergie électrique $\mathcal{E}_{e}$ emmagasinée dans le condensateur, l'énergie magnétique $\mathcal{E}_{m}$ emmagasinée dans la bobine et l'énergie totale $\mathcal{E}$ du circuit, tel que $\mathcal{E}=\mathcal{E}_{e}+\mathcal{E}_{m}$.

\begin{figure}[H]
	\centering
	\includegraphics[width=.7\textwidth]{22adf374-91fd-4329-b77c-a241db44a97c-5_604_1287_443_463}
	\caption{Figure (4)}
\end{figure}


2.5.1. Identifier, en justifiant la réponse, la courbe qui correspond à l'énergie magnétique $\mathcal{E}_{m}$.\\
2.5.2. Déterminer, entre les instants $t_{0}=0$ et $t_{1}=3 \mathrm{~ms}$, la variation $\Delta \mathcal{E}$ de l'énergie totale du circuit.

\section*{Exercice 3 ( $\mathbf{5}$ points) : Étude du mouvement d'un cycliste dans un circuit}
La course à bicyclette sur des circuits fermés est devenue un sport très populaire. Plusieurs compétitions s'organisent chaque année avec des circuits fermés qui comprennent des obstacles. Cet exercice vise l'étude du mouvement du centre d'inertie d'un système \{Cycliste - Bicyclette\} dans un circuit fermé de la région de l'Atlas (figure 1).\\
Au cours de sa participation à une course dont le circuit est représenté sur la figure (1), un cycliste parcourt une partie de ce circuit constituée d'un tronçon AB rectiligne horizontal, d'un tronçon BC curviligne qui s'ouvre sur une fosse de largeur L et d'un tronçon DE horizontal (figure 2).

\begin{figure}[H]
	\centering
	\includegraphics[width=.7\textwidth]{22adf374-91fd-4329-b77c-a241db44a97c-5_734_736_1722_242}
	\caption{Figure (1)}
\end{figure}


\begin{figure}[H]
	\centering
	\includegraphics[width=.7\textwidth]{22adf374-91fd-4329-b77c-a241db44a97c-5_501_966_1722_979}
	\caption{Figure (2)}
\end{figure}


Le mouvement sur le tronçon AB se fait avec des frottements modélisés par une force $\vec{f}$ constante de sens opposé au sens du vecteur vitesse. L'ensemble \{Cycliste - Bicyclette\} constitue un système de masse $m$ et de centre d'inertie $G$.

\section*{1. Mouvement du cycliste sur le tronçon $\mathbf{A B}$}
Le cycliste exerce entre A et B un effort modélisé par une force $\vec{F}$ horizontale supposée constante de même sens que le mouvement de $G$.\\
Le cycliste démarre sans vitesse initiale de la position A. Pour étudier le mouvement de $G$, on choisit le repère ( $A, \vec{i}$ ) lié à la Terre supposé Galiléen. À l'instant $t_{0}=0, x_{G}=x_{A}=0$.

\section*{Données :}
$m=70 \mathrm{~kg} \quad ; g=10 \mathrm{~m} . \mathrm{s}^{-2} \quad ; \quad F=180 N ; f=80 N ; \quad A B=60 \mathrm{~m}$\\
1.1. En appliquant la deuxième loi de Newton, montrer que l'expression de l'accélération du mouvement de G s'écrit : $a=\frac{F-f}{m}$.\\
1.2. Déterminer, en justifiant la réponse, la nature du mouvement de G .\\
1.3. Calculer la valeur de $t_{B}$, instant de passage de $G$ par $B$.\\
1.4. Déterminer la valeur de la vitesse $v_{B}$ de $G$ lors de son passage par $B$.\\
1.5. Déterminer l'intensité de la force $\vec{R}$ exercée par le plan sur le système au cours de son mouvement sur le tronçon $A B$.

\section*{2. Mouvement du cycliste durant la phase du saut}
Le cycliste quitte le tronçon BC en C avec une vitesse $\vec{v}_{0}$ qui fait un angle $\alpha$ avec le plan horizontal (voir figure 2- page 5/6).\\
Au cours du saut, le système \{Cycliste - Bicyclette\} n'est soumis qu'à son poids. On étudie le mouvement de $G$, dans un repère orthonormé ( $\mathrm{C}, \vec{i}^{\prime}, \vec{j}^{\prime}$ ) lié à la Terre supposé Galiléen. On choisit l'instant de passage de $G$ en C comme nouvelle origine des dates $t_{0}=0$.\\
Les équations horaires du mouvement de $G$ lors de la chute libre s'écrivent:

$$
x_{G}(t)=\left(v_{0} \cdot \cos \alpha\right) \cdot t \quad ; \quad y_{G}(t)=-\frac{1}{2} g \cdot t^{2}+\left(v_{0} \cdot \sin \alpha\right) \cdot t
$$

Au cours du mouvement, $G$ atteint le sommet de la trajectoire à l'instant $t_{s}=0,174 \mathrm{~s}$ et puis le système tombe sur le sol à l'instant $t_{P}=1 \mathrm{~s}$.

\section*{Données:}
$$
\alpha=10^{\circ} ; L=8 \mathrm{~m} ; g=10 \mathrm{~m} . \mathrm{s}^{-2}
$$

2.1. Montrer que $v_{0}=10 \mathrm{~m} . \mathrm{s}^{-1}$.\\
2.2. Le cycliste a-t-il dépassé la fosse ? justifier.\\
2.3. Déterminer les coordonnées du vecteur vitesse $\vec{v}_{P}$ de $G$ à l'instant $t_{P}$.


